\section{Results}
\label{sec:results}

\subsection{Headline results}
\label{sec:results:headline}

Table~\ref{tab:headline} reports the linear-probe LOSO numbers for our
default configuration (2.86\,M parameters; pretrained on subjects
51--100, 10{,}000 steps at $\lambda{=}1.0$; evaluated on the strictly
held-out subjects 1--20) on both EEGMMIDB binary tasks. The strongest and
\emph{only consistently significant} feature source on both tasks is the
simple \texttt{encoder\_mean}: 68.6\,\%\,$\pm$\,9.9 accuracy / 0.751 AUC on
rest-vs-activity ($+11.7$\,pp over the no-pretraining control;
$p{<}10^{-5}$) and 62.1\,\%\,$\pm$\,8.4 / 0.690 AUC on left-vs-right MI
($+10.9$\,pp; $p{=}4{\times}10^{-4}$). All $p$-values are exact paired
sign-flip permutation tests over the 20 LOSO folds, Benjamini--Hochberg
corrected (\cref{sec:results:significance}).

\begin{table*}[t]
\caption{Linear-probe LOSO accuracy and AUC on EEGMMIDB across 20
strictly held-out subjects (subjects~1--20; pretraining on 51--100).
\emph{pre}~=~pretrained checkpoint; \emph{rnd}~=~random-initialized
identical architecture. $\Delta$ is the pretrained$-$random accuracy gap;
$q$ is the Benjamini--Hochberg-corrected $p$-value from the paired
permutation test. Best per task in bold;
\textsuperscript{$\dagger$} marks gaps that survive correction
($q<0.05$).}
\label{tab:headline}
\centering
\small
\begin{tabular}{l l c c c c c}
\toprule
Task & Source & pre acc & pre AUC & rnd acc & $\Delta$ (pp) & $q$ \\
\midrule
left vs.\ right MI    & \textbf{encoder\_mean} & \textbf{0.621\,$\pm$\,0.084} & \textbf{0.690} & 0.512 & $\mathbf{+10.9}$\textsuperscript{$\dagger$} & 0.001 \\
                      & predictor\_mean        & 0.606\,$\pm$\,0.082          & 0.623          & 0.596 & $+1.0$                                       & 0.81  \\
                      & both\_mean             & 0.612\,$\pm$\,0.075          & 0.678          & 0.603 & $+0.9$                                       & 0.81  \\
\midrule
rest vs.\ activity    & \textbf{encoder\_mean} & \textbf{0.686\,$\pm$\,0.099} & \textbf{0.751} & 0.569 & $\mathbf{+11.7}$\textsuperscript{$\dagger$} & 1.5\textsc{e}-5 \\
                      & predictor\_mean        & 0.673\,$\pm$\,0.088          & 0.756          & 0.618 & $+5.4$\textsuperscript{$\dagger$}            & 0.002 \\
                      & both\_mean             & 0.679\,$\pm$\,0.100          & 0.757          & 0.638 & $+4.2$\textsuperscript{$\dagger$}            & 0.022 \\
\bottomrule
\end{tabular}
\end{table*}

A clear pattern emerges. \texttt{encoder\_mean}---the simplest source,
i.e.\ mean-pooled per-patch encoder embeddings with no temporal
context---delivers the largest pretraining benefit on both tasks and is
the only source whose gap over random is significant on the harder
left-vs-right task ($+10.9$\,pp vs.\ $+1.0$ and $+0.9$\,pp, both n.s., for
\texttt{predictor\_mean}/\texttt{both\_mean}). On rest-vs-activity all
three sources beat random, but the gap is again largest for the encoder.
We note that an earlier configuration that did not strictly hold out the
evaluation subjects from pretraining produced a markedly larger and more
predictor-favorable picture; under the leakage-free split reported here,
the predictor-derived sources lose most of their apparent advantage,
which we attribute to subject-level information having leaked through the
predictor's temporal aggregation
(\cref{sec:results:significance,sec:disc:predictor}).

\subsection{Scaling behaviour}
\label{sec:results:scaling}

We characterise the small-scale regime by training six configurations
spanning $\sim$$170\times$ in pretraining-data scale ($3$ to $109$
subjects) and $150\times$ in step count ($200$ to $30{,}000$ steps).
All operating points pretrain on subjects disjoint from the evaluation set
(eval $=$ subjects 1--20; pretraining drawn from 21--109; the full-corpus
points use the 89 subjects 21--109). \Cref{fig:scaling} plots best-source
AUC on rest-vs-activity against \emph{total window-presentations}, defined
as $N_\text{steps} \times \text{batch}$---the number of 4\,s EEG windows the
optimizer consumed (counting repeats), on a logarithmic axis.

\begin{figure}[t]
  \centering
  \IfFileExists{../runs/scaling_figure.png}
    {\includegraphics[width=\columnwidth]{scaling_figure.png}}
    {\fbox{\parbox{0.9\columnwidth}{\centering\vspace{2em}
      \texttt{scaling\_figure.png} — generate via
      \texttt{scripts/13\_scaling\_figure.py}\vspace{2em}}}}
  \caption{Compact-regime scaling (leakage-free; eval subjects 1--20 held
  out from all pretraining). Best-source rest-vs-activity LOSO AUC vs.\
  total window-presentations ($N_\text{steps}\times\text{batch}$, log axis).
  The base 2.86\,M model rises then plateaus near 0.76\,AUC; the
  7\,M-parameter model (diamond) sits \emph{above} the plateau at matched
  data and compute, indicating the base-model plateau reflects model
  capacity rather than a hard data limit.}
  \label{fig:scaling}
\end{figure}

Three observations frame the scaling picture. First, best-source AUC
improves steeply over the first two orders of magnitude of compute, from
0.710 at the 3-subject/200-step pilot to $\sim$0.76 by
$\sim$3$\times$10$^{5}$ window-presentations. Second, the base model then
\emph{plateaus} near 0.75--0.77\,AUC: the 50-subject/10k, 89-subject/10k,
and 89-subject/30k runs all land in this band, and pushing the 89-subject
corpus from 10k to 30k steps slightly \emph{lowers} AUC (0.768
$\rightarrow$ 0.750), a mild SSL-overfitting effect at fixed corpus. At
matched window-presentations, more subjects help (89-subject/10k $>$
50-subject/10k), consistent with data still mattering within the base
model's capacity. Third---and unlike the contaminated configuration we
initially examined---the 2.4$\times$ larger 7\,M model trained on the same
89-subject corpus reaches the highest AUC of all (0.774), sitting above the
base-model plateau. The base-model plateau therefore reflects a
\emph{capacity} limit at 2.86\,M parameters, not a hard data limit
(\cref{sec:results:capacity}).

\subsection{$\lambda$ ablation}
\label{sec:results:lambda}

We swept the SIGReg weight across
$\lambda \in \{0, 0.1, 0.3, 1.0, 3.0, 10.0\}$ at our default training
recipe (subjects 51--100, 10{,}000 steps, batch 64).
\texttt{encoder\_mean} rest-vs-activity accuracy is summarised in
Table~\ref{tab:lambda}.

\begin{table}[t]
\caption{SIGReg-weight ablation. \texttt{encoder\_mean} accuracy / AUC on
EEGMMIDB rest-vs-activity LOSO, 20 held-out subjects. $\Delta$ is over the
random-init baseline (0.569 acc). The $\lambda{=}0$ row is the SIGReg-off
control; it collapses to exactly chance.}
\label{tab:lambda}
\centering
\small
\begin{tabular}{c c c c}
\toprule
$\lambda$ & accuracy & AUC & $\Delta$ (pp) \\
\midrule
0.0   & 0.500 & 0.500 & $-6.8$ \\
0.1   & 0.649 & 0.711 & $+8.1$  \\
0.3   & 0.678 & 0.738 & $+11.1$  \\
\textbf{1.0} & \textbf{0.686} & \textbf{0.751} & \textbf{$+11.8$} \\
3.0   & 0.671 & 0.731 & $+10.3$  \\
10.0  & 0.619 & 0.661 & $+5.2$  \\
\bottomrule
\end{tabular}
\end{table}

Two findings are worth highlighting. First, $\lambda{=}1.0$ is a clean
optimum: a 10$\times$ decrease ($\lambda{=}0.1$) costs $\sim$4\,pp accuracy
and a 10$\times$ increase ($\lambda{=}10.0$) costs $\sim$7\,pp accuracy.
The downstream-accuracy curve is U-shaped with a sharp minimum, but the
training itself never destabilizes across the full 100$\times$ swing in
$\lambda$, validating SIGReg's reported stability claim from the original
LeJEPA paper.

Second---and more importantly---the $\lambda{=}0$ run \emph{collapses to
exactly chance accuracy} across both downstream tasks and all three
feature sources. This is the JEPA-collapse failure mode that SIGReg was
designed to prevent: without the distribution-matching regularizer, the
model finds the trivial solution of mapping every input to the same
embedding, and the linear probe receives no signal. The collapse is
total, not partial, which makes it the empirically-cleanest possible
validation of SIGReg's anti-collapse function. This complements the
theoretical analysis of \cite{lejepa2025}.

\subsection{Capacity test}
\label{sec:results:capacity}

To test whether the base-model plateau reflects a data limit or a capacity
limit, we trained a 2.4$\times$ larger variant ($7$\,M parameters; encoder
dim 256/mlp\_depth 3, predictor dim 256/depth 6) on the same leakage-free
89-subject corpus (subjects 21--109) for 30{,}000 steps. To separate
capacity from the SIGReg projection count $M$, we ran the larger model at
two projection counts ($M{=}1024$ and $M{=}256$) and repeated both the base
and large $M{=}1024$ configurations across three seeds (42--44). All
variants trained stably.

\begin{table}[t]
\caption{Capacity test: base (2.86\,M params) vs.\ large (7\,M params),
both trained leakage-free on the 89-subject corpus (subjects 21--109) for
30{,}000 steps at matched SIGReg projection count $M{=}1024$. EEGMMIDB LOSO
across 20 held-out subjects; accuracy as mean\,$\pm$\,std over three seeds
(42--44). $\Delta$ is large $-$ base. The larger model wins on five of six
cells and on both \texttt{encoder\_mean} (our primary source) rows.}
\label{tab:capacity}
\centering
\small
\begin{tabular}{l l c c c}
\toprule
Task & Source & base & large & $\Delta$ (pp) \\
\midrule
rva & \textbf{encoder\_mean}   & 0.678\,$\pm$\,0.008 & 0.692\,$\pm$\,0.011 & $+1.3$  \\
rva & predictor\_mean & 0.681\,$\pm$\,0.005 & 0.690\,$\pm$\,0.010 & $+0.9$  \\
rva & both\_mean      & 0.678\,$\pm$\,0.006 & 0.691\,$\pm$\,0.010 & $+1.3$  \\
\midrule
lr  & \textbf{encoder\_mean}   & 0.630\,$\pm$\,0.006 & 0.650\,$\pm$\,0.009 & $+2.0$  \\
lr  & predictor\_mean & 0.598\,$\pm$\,0.017 & 0.617\,$\pm$\,0.011 & $+2.0$  \\
lr  & both\_mean      & 0.642\,$\pm$\,0.006 & 0.641\,$\pm$\,0.007 & $-0.1$  \\
\bottomrule
\end{tabular}
\end{table}

At matched projection count ($M{=}1024$) the larger model outperforms the
base on five of six (task, source) cells (Table~\ref{tab:capacity}); on our
primary \texttt{encoder\_mean} source it wins on both tasks ($+2.0$\,pp on
left-vs-right, $+1.3$\,pp on rest-vs-activity). The margins are small but
remarkably stable: on left-vs-right \texttt{encoder\_mean} every large-model
seed (0.641--0.659) exceeds every base-model seed (0.623--0.636)---the two
seed distributions do not overlap---despite a per-seed spread under 1\,pp.
More capacity therefore helps at this data scale, and the base-model plateau
of \cref{sec:results:scaling} is a (mild) capacity limit rather than a hard
data limit. We nonetheless frame the effect as ``modestly better'' rather
than a scaling law: the $\sim$1--2\,pp margins, though seed-stable, remain
within the per-subject noise band.

A separate run isolates why an earlier, contaminated version of this
experiment (evaluation subjects present in pretraining; $M{=}256$ SIGReg
slices) had shown the \emph{opposite} result, with the larger model
appearing to regress. Holding architecture, data, and step count fixed and
varying only the projection count, the large model's left-vs-right
\texttt{encoder\_mean} accuracy collapses from 0.650 at $M{=}1024$ to 0.567
at $M{=}256$---an 8.3\,pp drop that pushes it \emph{below} the base model
(0.630) and reproduces the original ``large regresses'' artifact.
Rest-vs-activity, the easier task, is nearly unaffected (0.692 $\to$ 0.687).
The larger model's 256-dimensional embedding is under-regularized at
$M{=}256$: the Cram\'er--Wold projection count must scale with embedding
dimension, so the $M{=}256$ default that suffices for the 192-dimensional
base model is not safe for wider encoders. We flag this as a practical
pitfall when applying SIGReg-style isotropy regularizers across model
widths.

The practical takeaway is unchanged: the compact 2.86\,M model captures most
of the achievable accuracy---within $\sim$1--2\,pp of the 2.4$\times$ larger
model on \texttt{encoder\_mean}---while being markedly cheaper to store and
run, the relevant trade-off for edge deployment.

\subsection{Cross-dataset transfer}
\label{sec:results:transfer}

We evaluate cross-dataset transfer from EEGMMIDB (64 channels, motor
imagery) to BCI-IV-2a (22 channels, four-class motor imagery) under two
protocols: \emph{channel-padded} (the 64-channel pretrained encoder
applied at inference with the 42 missing channels zero-padded), and
\emph{channel-matched} (re-pretrain the same architecture on the 22-channel
intersection of EEGMMIDB before probing BCI-IV-2a). Results appear in
Table~\ref{tab:transfer}.

\begin{table}[t]
\caption{Cross-dataset transfer EEGMMIDB $\to$ BCI-IV-2a 4-class MI,
LOSO across 9 subjects. $\Delta$ is accuracy over the random-initialized
identical architecture (chance $=0.25$). The padded protocol applies the
64-channel pretrained encoder and zero-pads the 42 missing channels at
inference; the matched protocol re-pretrains on the 22-channel
intersection. Higher $\Delta$ is better; the padded protocol wins on every
source.}
\label{tab:transfer}
\centering
\small
\begin{tabular}{l c c c}
\toprule
Source & Padded $\Delta$ & Matched $\Delta$ & Padded $-$ Matched \\
\midrule
\textbf{encoder\_mean}   & $\mathbf{+4.9}$ & $+1.5$ & $+3.4$ \\
predictor\_mean & $+4.2$ & $+0.4$ & $+3.8$ \\
both\_mean      & $+2.0$ & $-1.5$ & $+3.5$ \\
\bottomrule
\end{tabular}
\end{table}

Two findings stand out. First, pretraining transfers across datasets, but
only modestly: the best configuration---\texttt{encoder\_mean} under the
padded protocol---reaches 0.314 accuracy / 0.583 macro-AUC on the 4-class
task ($+4.9$\,pp over the random-init control; chance $=0.25$). Absolute
accuracy stays well below the within-dataset numbers of
\cref{sec:results:headline}, consistent with our scope: we transfer to a
channel \emph{subset} of the pretraining montage, and a fully disjoint
layout would require the dynamic channel-mixing of \cite{laya2026}
(\cref{sec:discussion}).

Second, and more usefully for practitioners, the naive channel-padded
protocol \emph{outperforms} channel-matched re-pretraining on every feature
source (Table~\ref{tab:transfer}). On \texttt{encoder\_mean}, padding the
full 64-channel encoder beats re-pretraining on the 22-channel intersection
by 3.4\,pp ($+4.9$ vs.\ $+1.5$); the matched protocol even goes slightly
negative on \texttt{both\_mean}. This is counterintuitive---training the
encoder on the channel layout it will see at inference ``should'' help---but
points to a representational-redundancy effect: the 64-channel encoder,
trained on spatially redundant motor-cortex data, learns several spatial
views of each underlying signal and degrades gracefully to one of them when
the non-overlapping channels are zeroed, whereas the 22-channel encoder is
forced into a tighter, less redundant representation that transfers less
well to subjects with different spatial activation patterns. The practical
recommendation is therefore to pretrain at the richest available channel
topology and zero-pad at inference, rather than re-pretrain per downstream
montage.

\subsection{Per-subject variance}
\label{sec:results:perfold}

Figure~\ref{fig:perfold} shows the per-subject LOSO accuracy of the clean
\texttt{encoder\_mean} checkpoint. Two subjects fall below chance on
left-vs-right MI and one on rest-vs-activity---matching the
supervised-from-scratch baseline (two and one, respectively). Per-fold
standard deviation is comparable to or slightly above supervised
(0.084 vs.\ 0.077 on left-vs-right; 0.099 vs.\ 0.072 on rest-vs-activity).
SSL pretraining at this scale therefore neither reduces per-fold variance
nor eliminates below-chance subjects; its advantage over supervised
training is reuse across tasks (\cref{sec:results:supervised}), not
per-subject robustness.

\begin{figure}[t]
  \centering
  \IfFileExists{../runs/per_fold_figure.png}
    {\includegraphics[width=\columnwidth]{per_fold_figure.png}}
    {\fbox{\parbox{0.9\columnwidth}{\centering\vspace{2em}
      \texttt{per\_fold\_figure.png} — generate via
      \texttt{scripts/08\_per\_fold\_figure.py}\vspace{2em}}}}
  \caption{Per-subject LOSO accuracy of the clean \texttt{encoder\_mean}
  checkpoint on the two EEGMMIDB binary tasks (left: left-vs-right MI;
  right: rest-vs-activity), pretrained vs.\ random-init, with chance and
  mean reference lines.}
  \label{fig:perfold}
\end{figure}

\subsection{Supervised comparison}
\label{sec:results:supervised}

To anchor the SSL$+$probe results against an end-to-end alternative,
we trained the same encoder architecture from scratch with a linear
classification head on each LOSO fold (cross-entropy, 500 steps, batch
16, AdamW at $5{\times}10^{-4}$). Both pipelines use the
identical 1.08\,M-parameter encoder.

\begin{table}[t]
\caption{SSL$+$linear probe (\texttt{encoder\_mean}) vs.\
supervised-from-scratch on EEGMMIDB, LOSO across 20 held-out subjects,
accuracy$\,\pm\,$std / AUC. $\Delta$ is the SSL$-$supervised accuracy
gap; \textsuperscript{$\dagger$} marks $q<0.05$ (paired permutation,
BH-corrected).}
\label{tab:supervised}
\centering
\small
\begin{tabular}{l c c c}
\toprule
Task & Supervised & SSL$+$probe & $\Delta$ (pp) \\
\midrule
left vs.\ right    & 0.623\,$\pm$\,0.077 / 0.707 & 0.621\,$\pm$\,0.084 / 0.690 & $-0.2$ \\
rest vs.\ activity & 0.601\,$\pm$\,0.072 / 0.663 & 0.686\,$\pm$\,0.099 / 0.751 & $+8.4$\textsuperscript{$\dagger$} \\
\bottomrule
\end{tabular}
\end{table}

SSL$+$probe ties supervised-from-scratch on left-vs-right MI ($-0.2$\,pp;
not significant, $p{=}0.97$) and significantly exceeds it on
rest-vs-activity ($+8.4$\,pp; $p{<}10^{-3}$, $q{=}4\times10^{-4}$). The
per-fold standard deviation and below-chance counts are comparable across
the two pipelines (\cref{sec:results:perfold}), so the value of SSL here is
not lower variance but reuse: a single pretrained encoder serves both
downstream tasks (and any future task) at sub-second per-task probing cost,
whereas the supervised baseline retrains the full network per task.

Compute economy makes the comparison even sharper. The supervised
pipeline retrains the full network per LOSO fold and per task: $20$
folds $\times 2$ tasks $= 40$ training runs at $\sim$45 seconds each on
the M1. The SSL pipeline pretrains once (10 minutes on the RTX 5090),
extracts features once per task ($<$10 seconds), then fits
\texttt{LogisticRegression} per fold (sub-second). The pretrained
encoder is reused across both downstream tasks; the supervised model is
not.

\subsection{Statistical significance}
\label{sec:results:significance}

Because per-fold standard deviations are large relative to several reported
gaps, we test each headline difference for significance across the LOSO
folds rather than relying on mean differences alone. For each comparison we
compute the mean paired difference, a 95\,\% percentile bootstrap confidence
interval (10{,}000 resamples), an exact paired sign-flip permutation
$p$-value (full $2^{n}$ enumeration at $n{=}20$ folds), a Wilcoxon
signed-rank $p$-value (normal approximation with tie and continuity
correction) as a rank-based cross-check, and Cohen's $d_z$ effect size.
Across the family of accuracy comparisons we control the false-discovery
rate with the Benjamini--Hochberg procedure. The full procedure is
implemented dependency-free in \texttt{scripts/12\_significance.py} (with a
correctness self-test) and the table below is generated automatically from
the per-fold outputs.
After correction, both \texttt{encoder\_mean} gaps over random remain
strongly significant ($q{=}0.001$ on left-vs-right MI; $q{=}1.5\times
10^{-5}$ on rest-vs-activity). The \texttt{predictor\_mean} and
\texttt{both\_mean} gaps on left-vs-right MI do \emph{not} survive
correction ($q{>}0.8$), confirming that the apparent advantage of the
predictor-derived sources was a function of the (now-removed) subject
leakage rather than a robust effect.

\IfFileExists{../runs/significance_table.tex}
  {% Auto-generated by scripts/12_significance.py
\begin{table}[t]
\caption{Paired significance of the headline gaps (LOSO folds). $\Delta$ is mean paired difference in accuracy (pp); CI is a 95\% percentile bootstrap; $p$ is an exact paired sign-flip permutation test; $q$ is the Benjamini--Hochberg FDR-adjusted value across the accuracy family; $d_z$ is the paired effect size.}
\label{tab:significance}
\centering\small
\begin{tabular}{l l l c c c c}
\toprule
Comparison & Task & Source & $\Delta$ (pp) & 95\% CI (pp) & $p$ / $q$ & $d_z$ \\
\midrule
Pre $-$ Rand & L-vs-R MI & both & +0.9 & [-2.9, +4.4] & 0.690 / 0.808 & +0.10 \\
Pre $-$ Rand & L-vs-R MI & enc & +10.9 & [+6.2, +16.0] & 0.000 / 0.001 & +0.94 \\
Pre $-$ Rand & L-vs-R MI & pred & +1.0 & [-3.7, +5.1] & 0.707 / 0.808 & +0.10 \\
Pre $-$ Rand & rest-vs-act & both & +4.2 & [+1.3, +6.9] & 0.014 / 0.022 & +0.63 \\
Pre $-$ Rand & rest-vs-act & enc & +11.7 & [+9.1, +14.6] & 0.000 / 0.000 & +1.81 \\
Pre $-$ Rand & rest-vs-act & pred & +5.4 & [+2.9, +7.8] & 0.001 / 0.002 & +0.95 \\
\midrule
SSL $-$ Sup & L-vs-R MI & enc & -0.2 & [-5.0, +4.7] & 0.965 / 0.965 & -0.02 \\
SSL $-$ Sup & rest-vs-act & enc & +8.4 & [+5.3, +11.8] & 0.000 / 0.000 & +1.10 \\
\bottomrule
\end{tabular}
\end{table}
}
  {}
